\section{Introduction}\label{sec:intro}

A surface code logical qubit is protected only as long as its syndrome is decoded fast and well
\cite{dennis2002topological,fowler2012surface}. Experiments that operate below the surface code threshold
\cite{google2023scaling,google2025below} have made decoding speed and accuracy questions of system design, and
minimum weight perfect matching, in the Sparse Blossom implementation of PyMatching
\cite{higgott2022pymatching,higgott2025sparse}, is the common reference for both. Neural pre-decoders are one proposal
to improve on it: a small convolutional network looks at a window of syndrome rounds, removes the errors it is sure
of, and leaves a sparser syndrome to the matching decoder
\cite{chamberland2018deep,chamberland2026predecoder}. NVIDIA released two such networks with pretrained weights and
evaluation code \cite{chamberland2026predecoder,nvidia_fast_card,nvidia_ising_code}. In earlier work we implemented them
as an integer kernel on Alveo U55C cards and measured their throughput on four cards \cite{ali2026nodes}.

Those networks were trained, and their published accuracy measured, on circuit noise of one strength at every
location. A device is different. On the IBM devices studied here the two qubit gate error of the worst usable coupler
is \WorstCouplerMin{} to \WorstCouplerMax{} times that of the median one, the median readout error is
\ReadoutOverGateMin{} to \ReadoutOverGateMax{} times the median two qubit gate error, and data qubits idle through
microseconds of ancilla readout with their own $T_1$ and $T_2$. Whether a network trained on uniform noise keeps its
advantage under such noise, and against which matching baseline that advantage should be judged, is the subject of
this paper. A decoder that knows each qubit's rates can also weight its matching graph accordingly
\cite{sivak2024priors}, which raises a second question: how much of any gain comes from the network and how much is
available to matching alone.

We answer both questions with device calibration data, without access to the devices. IBM publishes calibration
snapshots of its devices with qiskit-ibm-runtime \cite{qiskit_ibm_runtime}; we export eight of them to CSV files in the
layout of IBM's own calibration downloads and map every fault location of NVIDIA's memory circuit to the rate of the
qubit or coupler it is assigned to. The pre-decoder then runs unchanged, in floating point and as the integer model
that the FPGA kernel computes, and on one U55C card.

The contributions are:
\begin{itemize}
\item A noise model that keeps the fault locations of NVIDIA's circuit, so the network sees the structure it was
trained on, and sets every location from a calibration CSV: gate errors, readout and preparation errors, and $T_1$ and
$T_2$ decay over the measured durations of gate, readout and reset layers (Section~\ref{sec:methods}).
\item An evaluation over \NumPoints{} points (eight devices, three placements, three error scales, distances 3 and 5)
with paired, shot by shot comparisons against plain matching, matching with device median weights and correlated
matching \cite{fowler2013correlated,paler2023pipelined} (Section~\ref{sec:results}).
\item The finding that correlated matching, with no network, beats the pre-decoder pipeline at every distance 5
point, that the network loses its advantage on strongly uneven qubits, and that calibrated matching weights are worth
up to \MedWeightsMax{} times in logical error on their own.
\item Integer inference of the network on one U55C card, bit exact on every run and identical to the CPU integer
model in every logical failure wherever both were run, with a direct throughput comparison against a \CpuThreads{} thread CPU.
\item A public repository with the calibration files, the experiments and a teaching example that runs codes of
distance 1, 2, 3 and 5 from any calibration CSV on a laptop.
\end{itemize}
